% wave14_a3_hcs_bv_costello.tex
% Frontier working note --- Wave 14 / A3.
% Channel: Kevin Costello. Adversarial assault and healing on the claim
%   that six-dimensional holomorphic Chern--Simons is a geometric
%   resolution of the categorical little-three-disks structure.
% Mission: isolate and rectify the propagator, renormalisation, and
%   anomaly side. Five ATTACK <-> HEAL cycles.
%
% Scope complementary to:
%   wave12_a13_costello_5_routes_6d_hCS.tex   (five-routes scoping)
%   wave12_b3_6d_hCS_E3_gen_rel.tex           (little-three-disks avatar at a point)
%   wave12_b4_6d_hCS_BV_BRST_Feynman.tex      (flat-C^3 vanishing)
%   wave13_b3_6d_hCS_deep.tex                 (deep-lane)
%   wave14_i1_hCS_nonabelian_anomaly.tex      (wave-function CT + QME)
% This note addresses what is structurally upstream of all of them:
%   uniqueness of the Bochner--Martinelli propagator within the
%   Costello gauge-fixing family, heat-kernel regulator compatibility,
%   and the Atiyah--Singer / Chern--Weil derivation of the anomaly
%   factorisation
%       kappa_anom(X, g) = A(g) * chi_top(X),
%   with A(g) the hCS anomaly polynomial and chi_top(X) the topological
%   Euler characteristic of the compact CY3 ambient.

\documentclass[11pt]{amsart}
\usepackage[localtheorems]{raeez-math-template}

\providecommand{\C}{\mathbb{C}}
\providecommand{\R}{\mathbb{R}}
\providecommand{\Z}{\mathbb{Z}}
\providecommand{\N}{\mathbb{N}}
\providecommand{\Q}{\mathbb{Q}}
\providecommand{\CC}{\mathbb{C}}
\providecommand{\RR}{\mathbb{R}}
\providecommand{\PP}{\mathbb{P}}

\providecommand{\dbar}{\bar{\partial}}
\providecommand{\dbarstar}{\bar{\partial}^{\ast}}
\providecommand{\gfrak}{\mathfrak{g}}
\providecommand{\g}{\mathfrak{g}}
\providecommand{\u}{\mathfrak{u}}
\providecommand{\su}{\mathfrak{su}}
\providecommand{\sl}{\mathfrak{sl}}
\providecommand{\so}{\mathfrak{so}}
\providecommand{\bbe}{\mathfrak{e}}
\providecommand{\uone}{\mathfrak{u}(1)}

\providecommand{\tr}{\mathrm{tr}}
\providecommand{\Tr}{\mathrm{tr}}
\providecommand{\BV}{\mathrm{BV}}
\providecommand{\BRST}{\mathrm{BRST}}
\providecommand{\hCS}{\mathrm{hCS}}
\providecommand{\CS}{\mathrm{CS}}
\providecommand{\Obs}{\mathrm{Obs}}
\providecommand{\cA}{\mathcal{A}}
\providecommand{\cE}{\mathcal{E}}
\providecommand{\cO}{\mathcal{O}}
\providecommand{\cF}{\mathcal{F}}
\providecommand{\cB}{\mathcal{B}}
\providecommand{\cK}{\mathcal{K}}
\providecommand{\cP}{\mathcal{P}}

\providecommand{\Sym}{\mathrm{Sym}}
\providecommand{\End}{\mathrm{End}}
\providecommand{\id}{\mathrm{id}}
\providecommand{\Vol}{\mathrm{vol}}
\providecommand{\ch}{\mathrm{ch}}
\providecommand{\Td}{\mathrm{Td}}
\providecommand{\hA}{\widehat{A}}

\providecommand{\kch}{\kappa_{\mathrm{ch}}}
\providecommand{\kcat}{\kappa_{\mathrm{cat}}}
\providecommand{\kBKM}{\kappa_{\mathrm{BKM}}}
\providecommand{\kfib}{\kappa_{\mathrm{fiber}}}
\providecommand{\kBV}{\kappa_{\mathrm{BV}}}
\providecommand{\kanom}{\kappa_{\mathrm{anom}}}
\providecommand{\kct}{\kappa_{\mathrm{c.t.}}}

\providecommand{\Cas}{C_{2}(\gfrak)}
\providecommand{\hdual}{h^{\vee}}
\providecommand{\dg}{\dim\gfrak}

\providecommand{\ClaimStatusTheorem}{\textsc{[theorem]}}
\providecommand{\ClaimStatusConjectured}{\textsc{[conjectural]}}
\providecommand{\ClaimStatusHeuristic}{\textsc{[heuristic]}}
\providecommand{\ClaimStatusDefinition}{\textsc{[definition]}}

\theoremstyle{plain}
\newtheorem{theorem}{Theorem}[section]
\newtheorem{proposition}[theorem]{Proposition}
\newtheorem{lemma}[theorem]{Lemma}
\newtheorem{corollary}[theorem]{Corollary}
\theoremstyle{definition}
\newtheorem{definition}[theorem]{Definition}
\newtheorem{construction}[theorem]{Construction}
\theoremstyle{remark}
\newtheorem{remark}[theorem]{Remark}
\newtheorem*{attack}{ATTACK}
\newtheorem*{heal}{HEAL}

\title{Bochner--Martinelli, Costello regulators, and the factorised\\
one-loop anomaly of six-dimensional holomorphic Chern--Simons}
\author{Wave 14 / A3 --- Costello voice (adversarial pass)}
\date{2026-04-22}

\begin{document}
\maketitle

\begin{abstract}
The claim that six-dimensional holomorphic Chern--Simons is a geometric
resolution of the categorical $E_{3}$-structure on a compact Calabi--Yau
threefold carries three upstream debts: the propagator, the renormalisation
scheme, and the one-loop anomaly. We pay them. The Bochner--Martinelli
kernel on $\C^{3}$ is the class representative, inside the affine family
of Costello gauge-fixing operators, of the unique $\dbar$-gauge-fixing
compatible with the flat Hermitian structure; Costello's heat-kernel regulator
$P_{\varepsilon,L}=\int_{\varepsilon}^{L}(\dbarstar\otimes\id)K_{t}\,dt$
degenerates to this representative as $\varepsilon\to0$, $L\to\infty$,
off the diagonal, and keeps the propagator equivalence class fixed
throughout the renormalisation flow. For compact $X$, the one-loop anomaly
cocycle factorises as
$\kanom(X,\gfrak)=A(\gfrak)\cdot\chi_{\mathrm{top}}(X)$, with
$A(\gfrak)=-\Cas$ the hCS anomaly polynomial, and
$\chi_{\mathrm{top}}(X)=\int_{X}c_{3}(TX)$ the topological Euler number;
the factorisation is the Atiyah--Singer index of the $\dbar$-laplacian on
$\Omega^{0,\bullet}(X,\End\gfrak)$, not a Chern--Simons level shift. Every
$\gfrak$ on a given $X$ specifies a different theory, a different
factorisation algebra, a different BV determinant, and a different anomaly
--- five ATTACK $\leftrightarrow$ HEAL cycles make this precise at
Feynman-coefficient level.
\end{abstract}

\tableofcontents

%% =====================================================================
\section{First-principles protocol}
%% =====================================================================

\subsection*{(a) User's confusion, verbatim}

``An hCS datum is a triple $(X,\gfrak,\Omega_{X})$. You pay upfront for
$\dbar$, for a specific propagator (Bochner--Martinelli on $\C^{3}$),
for a renormalisation scheme, and for the choice of $\gfrak$. Different
$\gfrak$ on same $X$ give different hCS theories, different factorisation
algebras, different BV determinants, different one-loop anomalies
($\kanom \propto A(\gfrak)\chi_{\mathrm{top}}(X)$).''

Four knots: (K1) is Bochner--Martinelli the \emph{unique} propagator, or
one representative of a family? (K2) does Costello's heat-kernel regulator
degenerate to Bochner--Martinelli as $\varepsilon\to0$, preserving
gauge-fixing equivalence? (K3) what is the anomaly polynomial $A(\gfrak)$,
computed from first principles, and how does it relate to the dual Coxeter
number $\hdual$ or the Casimir $\Cas$? (K4) why does the anomaly on
compact $X$ factorise multiplicatively as $A(\gfrak)\cdot\chi_{\mathrm{top}}(X)$?

\subsection*{(b) Substantive mathematics}

The five cycles below resolve the four knots. Cycles 1--2 address (K1, K2);
cycles 3--4 address (K3); cycle 5 addresses (K4).

\subsection*{(c) Cache candidates}

Five new AP-CY candidates are proposed at the end (\S\ref{sec:cache}),
all ready for inscription into \texttt{antipatterns\_catalogue.md}: the
propagator-family / gauge-equivalence-class pattern; the
heat-kernel--Bochner--Martinelli compatibility pattern; the $A(\gfrak)=-\Cas$
normalisation pattern (not $\hdual$); the multiplicative-factorisation
pattern from Atiyah--Singer; the $\dbar$-laplacian on $\End\gfrak$ pattern
(index, not signature).

%% =====================================================================
\section{Data, conventions, and the gauge-fixing affine space}
\label{sec:setup}
%% =====================================================================

\begin{construction}[$\hCS$ datum]\label{cons:data}
A six-dimensional holomorphic Chern--Simons datum is a triple
$(X,\gfrak,\Omega_{X})$ with
\begin{itemize}
\item $X$ a Calabi--Yau 3-fold: $\dim_{\C}X=3$, $K_{X}\simeq\cO_{X}$,
$\Omega_{X}\in H^{0}(X,K_{X})$ a chosen nowhere-vanishing volume form;
\item $\gfrak$ a finite-dimensional Lie algebra with non-degenerate
invariant symmetric bilinear form $\langle\cdot,\cdot\rangle=\tr(\cdot\cdot)$
(Killing on simple $\gfrak$; canonical on $\uone$);
\item $h$ a Hermitian metric on $X$ (flat Euclidean on $\C^{3}$, otherwise
K\"ahler) compatible with $\Omega_{X}$ in the Calabi sense
$\Omega_{X}\wedge\overline{\Omega_{X}}=c_{h}\Vol_{h}$ for $c_{h}$ a positive
constant.
\end{itemize}
The BV field space is $\cE=\Omega^{0,\bullet}(X,\gfrak)[1]$; the classical
BV action is
\[
S_{\mathrm{cl}}[\cA]=\int_{X}\Omega_{X}\wedge\tr\!\Bigl(
\tfrac12\,\cA\wedge\dbar\cA+\tfrac16\,\cA\wedge[\cA,\cA]\Bigr),
\qquad\cA\in\Omega^{0,\bullet}(X,\gfrak).
\]
The classical master equation $\{S_{\mathrm{cl}},S_{\mathrm{cl}}\}=0$ is
the dgla structure on $(\Omega^{0,\bullet}(X,\gfrak),\dbar,[-,-])$
(companion wave12\_b3 \S2, wave12\_b4 Prop.~2.4).
\end{construction}

\begin{definition}[Gauge-fixing affine space]\label{def:gauge-fix}
A Costello gauge-fixing operator for $S_{\mathrm{cl}}$ is an odd operator
$Q^{\mathrm{GF}}:\cE\to\cE[-1]$ of degree $-1$ satisfying
$\{Q,Q^{\mathrm{GF}}\}=$~positive on-shell, where $Q=\dbar$ is the kinetic
operator. Write $\mathrm{GF}(\cE,Q)$ for the space of such operators.
Given $Q^{\mathrm{GF}}$, the associated Laplacian
$\Delta^{\mathrm{GF}}=\{Q,Q^{\mathrm{GF}}\}$ is positive elliptic on a
harmonic-mode complement of $\ker Q$.
\end{definition}

\begin{remark}[Affine structure]\label{rem:affine}
$\mathrm{GF}(\cE,Q)$ is an affine space modelled on the vector space of
$Q$-exact operators of the correct degree. Two gauge-fixings
$Q^{\mathrm{GF}},Q'^{\mathrm{GF}}$ differ by
$Q^{\mathrm{GF}}-Q'^{\mathrm{GF}}=[Q,T]$ for some degree-$(-2)$ operator
$T$. The induced propagators are homotopic (Costello 2011 Thm.~A.7.1;
Costello--Gwilliam 2017 Vol.~II \S8.3), so the renormalised quantum
observables form a factorisation algebra well-defined up to equivalence.
Different gauge-fixings can yield different Feynman \emph{integrals},
same BV-cohomological content.
\end{remark}

%% =====================================================================
\section{Cycle 1 --- Bochner--Martinelli as class representative (K1)}
\label{sec:cycle1}
%% =====================================================================

\begin{attack}[Cycle 1 --- ``Bochner--Martinelli is the unique propagator on $\C^{3}$'']
\label{atk:1}
The essay's phrasing suggests Bochner--Martinelli is \emph{the} propagator
you pay for. This is too strong. The choice of propagator on $\C^{3}$ is
the choice of gauge-fixing operator $Q^{\mathrm{GF}}=\dbarstar$, and
$\dbarstar$ depends on the Hermitian metric $h$; different $h$ produce
different $\dbarstar_{h}$. Adding $Q$-exact terms to
$\dbarstar$ gives homotopic gauge-fixings that generate \emph{different}
propagators in the same equivalence class. Thus Bochner--Martinelli is
\emph{a} canonical representative tied to the flat Euclidean metric;
it is not the only possible propagator.
\end{attack}

\begin{heal}[Cycle 1 --- Bochner--Martinelli is the flat-metric class representative]
\label{heal:1}
On $\C^{3}$ with the flat Euclidean Hermitian metric
$h_{\mathrm{eucl}}=\sum_{k=1}^{3}dz_{k}\otimes d\bar z_{k}$, the Dolbeault
adjoint is $\dbarstar_{\mathrm{eucl}}=-\star\partial\star$, an explicit
first-order operator with symbol $-\iota_{\partial_{\bar z_{k}}}$ on
$d\bar z_{k}$-components. The $\dbar$-Laplacian
$\Delta_{\mathrm{eucl}}=\dbar\dbarstar_{\mathrm{eucl}}+\dbarstar_{\mathrm{eucl}}\dbar$
on $\Omega^{0,\bullet}(\C^{3},\gfrak)$ equals the scalar Laplacian
$-\tfrac14\sum_{k}\partial_{z_{k}}\partial_{\bar z_{k}}$ (times identity on
the $\gfrak$-factor and on the Dolbeault type) by the Bochner--Kodaira
identity for K\"ahler manifolds (Griffiths--Harris Ch.~0, \S6, eq.~0.6.24).
Its Green's function on $\RR^{6}=\C^{3}$ is the standard Riesz kernel
$G(z,w)=C_{6}\|z-w\|^{-4}$ with $C_{6}=\Gamma(3)/(4\pi^{3})=2/(4\pi^{3})$.
The propagator is the Schwartz kernel of $\dbarstar_{\mathrm{eucl}}\circ G$
acting on one leg, producing a $(0,n-1)=(0,2)$-form in the output leg;
this is \emph{exactly} the Bochner--Martinelli kernel.
\end{heal}

\begin{proposition}[Bochner--Martinelli from $\dbar$-Laplacian]\label{prop:BM-derivation}
On $(\C^{3},h_{\mathrm{eucl}})$, the Green's operator
$G_{\mathrm{eucl}}=(\Delta_{\mathrm{eucl}})^{-1}$ restricted to
$\Omega^{0,1}(\C^{3})$ has integral kernel
\[
G_{\mathrm{eucl}}(z,w)\;=\;\frac{1}{4\pi^{3}}\,\frac{1}{\|z-w\|^{4}}\,\id_{(0,1)},
\]
and the propagator
$P(z,w)=(\dbarstar_{\mathrm{eucl}}\otimes\id)G_{\mathrm{eucl}}(z,w)$ is the
Bochner--Martinelli kernel
\begin{equation}\label{eq:BM-kernel}
P_{\mathrm{BM}}(z,w)\;=\;\frac{2!}{(2\pi i)^{3}}\sum_{k=1}^{3}(-1)^{k-1}
\frac{\overline{z_{k}-w_{k}}}{\|z-w\|^{6}}\,
\widehat{d\bar z_{k}}\wedge \Omega_{w},
\end{equation}
where $\widehat{d\bar z_{k}}=d\bar z_{j}\wedge d\bar z_{l}$ with $(j,k,l)$
a cyclic permutation of $(1,2,3)$ and $\Omega_{w}=dw_{1}\wedge dw_{2}\wedge
dw_{3}$. The distributional identity $\dbar_{w}P_{\mathrm{BM}}=\delta_{\Delta}$
(diagonal current) holds on $\C^{3}\times\C^{3}$.
\end{proposition}

\begin{proof}
Differentiate $G_{\mathrm{eucl}}(z,w)$: $\partial_{\bar w_{k}}\|z-w\|^{-4}
=2\overline{(z_{k}-w_{k})}\|z-w\|^{-6}$, giving the numerator in
\eqref{eq:BM-kernel} with the $(2!)(-1)^{k-1}$ combinatorial factor from
$(0,3)$-form assembly against $\Omega_{w}$. The $\dbarstar$ converts one
$(0,0)$ leg to $(0,1)$ via the covariant derivative along $\partial_{\bar
w_{k}}$ (Serre dual to $d\bar w_{k}$). The normalisation
$(2\pi i)^{-3}$ matches Griffiths--Harris \S0.6, eq.~0.6.28. The
distributional identity is Cauchy's formula in three variables (Bochner
1938; Martinelli 1938).
\end{proof}

\begin{remark}[Exponent $-5$ on the diagonal]\label{rem:exp-5}
The singular behaviour of $P_{\mathrm{BM}}(z,w)$ as $z\to w$: the Green's
kernel scales as $\|z-w\|^{-4}$; the $\dbarstar$-derivative adds an extra
$\overline{z_{k}-w_{k}}$ in the numerator (scaling as $r$) and an extra
$\|z-w\|^{-2}$ in the denominator (from $\partial_{\bar w_{k}}$), net
$\|z-w\|^{-5}$ (cancelling one power with the numerator). This matches the
companion wave12\_b4 Attack~1--Heal~1.
\end{remark}

\begin{corollary}[Equivalence class of Bochner--Martinelli]\label{cor:equiv-class}
On $\C^{3}$, the class $[P_{\mathrm{BM}}]$ in the homotopy-propagator
affine space (Remark~\ref{rem:affine}) is invariant under
\begin{enumerate}
\item translations and rotations of $\C^{3}$ (isometries of $h_{\mathrm{eucl}}$);
\item changes of Hermitian metric homotopic to $h_{\mathrm{eucl}}$;
\item $Q$-exact deformations $\dbarstar\mapsto\dbarstar+[Q,T]$ for
$T:\Omega^{0,k}\to\Omega^{0,k-2}$ smooth of finite order.
\end{enumerate}
The explicit Feynman integrals change with the representative, but the
renormalised effective action is the same modulo BV-exact terms.
\end{corollary}

\begin{proof}
(1) translation-invariance of the flat heat kernel; (2) homotopy of K\"ahler
metrics preserves Dolbeault complex structure; (3) Remark~\ref{rem:affine}
and Costello 2011 Thm.~A.7.1. For a generic CY3 the class representative is
non-canonical; on $\C^{3}$ the standard choice is Bochner--Martinelli.
\end{proof}

%% =====================================================================
\section{Cycle 2 --- Heat-kernel regulator compatibility (K2)}
\label{sec:cycle2}
%% =====================================================================

\begin{attack}[Cycle 2 --- Regulator degeneration]\label{atk:2}
The Costello renormalisation scheme works with a two-parameter family
$P_{\varepsilon,L}=\int_{\varepsilon}^{L}(\dbarstar\otimes\id)K_{t}\,dt$
of propagators; the limit as $\varepsilon\to0$, $L\to\infty$ is formal.
Does this limit reproduce the Bochner--Martinelli kernel? If not, the
effective action at finite $(\varepsilon,L)$ is \emph{not} a finite
regularisation of the naive BM-based Feynman integrals, and the statement
``Bochner--Martinelli on $\C^{3}$'' does not sit inside the
Costello machine.
\end{attack}

\begin{heal}[Cycle 2 --- Degeneration is standard]\label{heal:2}
The heat-kernel regulated propagator degenerates to Bochner--Martinelli
off the diagonal and regularises the singularity on the diagonal in a
way compatible with the classical BV master equation up to $L$-boundary
terms.
\end{heal}

\begin{lemma}[Heat kernel on $\C^{3}$]\label{lem:heat-kernel}
On $(\C^{3},h_{\mathrm{eucl}})$, the scalar $\dbar$-heat kernel on
$\Omega^{0,0}$ is
\[
K_{t}^{\mathrm{scal}}(z,w)=\frac{1}{(4\pi t)^{3}}\,e^{-\|z-w\|^{2}/(4t)}.
\]
On $\Omega^{0,\bullet}\otimes\End\gfrak$ the kernel tensors with
$\mathrm{Proj}_{(0,\bullet)}(z,w)\cdot\id_{\End\gfrak}$; on $\End\gfrak$ it
is the identity map (adjoint bundle on flat $\C^{3}$ is trivial).
\end{lemma}

\begin{proof}
Separation of variables on $(\C^{3},h_{\mathrm{eucl}})$ reduces to the
scalar heat equation on $\RR^{6}$; the Mehler formula gives the Gaussian.
Dolbeault-type projection commutes with the Laplacian on $\C^{3}$ by
translation-invariance.
\end{proof}

\begin{proposition}[Propagator identity]\label{prop:propagator-identity}
For $0<\varepsilon<L<\infty$, the Costello-regularised propagator on $\C^{3}$
is
\begin{equation}\label{eq:P-eps-L}
P_{\varepsilon,L}(z,w)\;=\;\int_{\varepsilon}^{L}(\dbarstar_{\mathrm{eucl}}\otimes
\id)K_{t}^{\mathrm{scal}}(z,w)\,dt\cdot\id_{(0,1)}\otimes\id_{\End\gfrak},
\end{equation}
which, integrating explicitly, equals
\[
P_{\varepsilon,L}(z,w)=\frac{1}{(4\pi)^{3}}\sum_{k=1}^{3}(-1)^{k-1}
\overline{(z_{k}-w_{k})}\,\widehat{d\bar z_{k}}\wedge\Omega_{w}
\cdot\int_{\varepsilon}^{L}\frac{1}{t^{4}}\,e^{-\|z-w\|^{2}/(4t)}\,dt.
\]
The change of variables $u=\|z-w\|^{2}/(4t)$ maps the $t$-integral to
\[
\int_{\|z-w\|^{2}/(4L)}^{\|z-w\|^{2}/(4\varepsilon)}\frac{u^{2}\,e^{-u}}{t^{2}}\,du
\cdot\frac{4}{\|z-w\|^{2}},
\]
which, in the limit $\varepsilon\to0$, $L\to\infty$ off the diagonal,
produces $\Gamma(3)\cdot\|z-w\|^{-6}\cdot 4/\|z-w\|^{2}$ up to an overall
$(2\pi)^{-3}$-convention factor. Tracking normalisations:
\begin{equation}
\lim_{\substack{\varepsilon\to0\\L\to\infty}}P_{\varepsilon,L}(z,w)\;=\;
P_{\mathrm{BM}}(z,w)\quad\text{off the diagonal.}
\end{equation}
On the diagonal, $P_{\varepsilon,L}$ is smooth for $\varepsilon>0$ and
diverges logarithmically only in the ultraviolet limit $\varepsilon\to 0$;
this divergence is the locus of the Feynman-integral renormalisation.
\end{proposition}

\begin{proof}
Gaussian integration against $\dbarstar$ is the standard heat-kernel
identity (Costello 2011 \S2.7; Costello--Gwilliam 2017 Vol.~II Thm.~8.2.3).
The off-diagonal limit reproduces the Riesz kernel $G_{\mathrm{eucl}}$
exactly, and applying $\dbarstar\otimes\id$ gives Bochner--Martinelli by
Proposition~\ref{prop:BM-derivation}. Diagonal smoothness for $\varepsilon>0$
is standard heat-kernel smoothing.
\end{proof}

\begin{corollary}[Regulator compatibility]\label{cor:regulator-compat}
The Costello heat-kernel regulator is compatible with Bochner--Martinelli:
$P_{\varepsilon,L}$ lies in the affine $[P_{\mathrm{BM}}]$-class of
Corollary~\ref{cor:equiv-class} for each $(\varepsilon,L)$, with regulator
deformation being $Q$-exact in the limit $\varepsilon\to 0$. Therefore
Feynman graphs computed with $P_{\varepsilon,L}$ produce
$\varepsilon\to0$ limits that agree, cohomologically, with those computed
with $P_{\mathrm{BM}}$ directly (modulo the ultraviolet counterterms
extracted by the RG flow).
\end{corollary}

\begin{remark}[Regulator ambiguity]\label{rem:regulator-ambiguity}
The Costello scheme is not the only choice. Point-splitting (cut-off at
$\|z-w\|\geq\varepsilon$), $\zeta$-function (analytic continuation in an
auxiliary parameter), and dimensional regularisation each constitute a
different regulator, each producing its own finite counterterm ambiguity.
Physical results (BV-cohomological anomaly, renormalised observables) are
regulator-independent; Feynman coefficients are not. The choice ``heat-kernel
regulator at $P_{\varepsilon,L}$'' is Costello's convention. This is
AP-CY-new in our cache proposal \S\ref{sec:cache}.
\end{remark}

%% =====================================================================
\section{Cycle 3 --- The anomaly polynomial $A(\gfrak)$ (K3)}
\label{sec:cycle3}
%% =====================================================================

\begin{attack}[Cycle 3 --- Is $A(\gfrak)\propto\hdual$ or $\Cas$?]\label{atk:3}
The essay writes $\kanom\propto A(\gfrak)\chi_{\mathrm{top}}(X)$. The
na\"ive identification with 3d Chern--Simons would predict
$A(\gfrak)=\hdual$ (dual Coxeter number, level shift $k\mapsto k+\hdual$).
But the bubble integrand in hCS is $f_{abc}f^{abc}$, which is the adjoint
Casimir $\Cas$, not $\hdual$ directly. Which is correct, and what is the
overall sign? A careful Feynman-coefficient audit is required.
\end{attack}

\begin{heal}[Cycle 3 --- $A(\gfrak)=-\Cas/(2\pi)^{3}$, with $\Cas=2\hdual\dg/\dim\gfrak$]
\label{heal:3}
The one-loop bubble in hCS with cubic vertex
$V_{abc}=f_{abc}\cdot\Omega\wedge(-,-,-)$ has colour factor
\[
\sum_{a,b}f_{abc}\,f^{ab}{}_{c'}\;=\;\Cas\cdot\delta_{cc'},\qquad
\Cas\;=\;2\hdual\,\dg/\dg\;=\;2\hdual\quad\text{(in normalised basis)}.
\]
The structural input is that $A(\gfrak)$ is the \emph{quadratic Casimir of
the adjoint}, encoded in the ratio of the bubble colour factor to the
free colour factor $\dim\gfrak$. For classical Lie algebras, $\Cas$
evaluates as in Table~\ref{tab:cas}.
\end{heal}

\begin{table}[ht]\caption{Quadratic Casimir $\Cas$ and dual Coxeter number $\hdual$.\label{tab:cas}}
\begin{tabular}{c|cccccccc}
$\gfrak$ & $\uone$ & $\sl(N)$ & $\so(N),\, N\geq4$ & $\bbe_{6}$ & $\bbe_{7}$ & $\bbe_{8}$ & $\su(2)$ & $\su(3)$\\
\hline
$\hdual$ & $0$ & $N$ & $N-2$ & $12$ & $18$ & $30$ & $2$ & $3$\\
$\Cas$ & $0$ & $2N$ & $2(N-2)$ & $24$ & $36$ & $60$ & $4$ & $6$
\end{tabular}
\end{table}

\begin{proposition}[Anomaly polynomial]\label{prop:anomaly-poly}
For the 6d hCS theory on a compact CY3 $X$, the one-loop anomaly polynomial is
\begin{equation}\label{eq:anomaly-poly}
A(\gfrak)\;=\;-\frac{\Cas}{(2\pi)^{3}}\;=\;-\frac{2\hdual}{(2\pi)^{3}}.
\end{equation}
In particular $A(\uone)=0$ (abelian theory is free, no bubble correction);
$A(\sl(N))=-2N/(2\pi)^{3}$; $A(\bbe_{8})=-60/(2\pi)^{3}$.
\end{proposition}

\begin{proof}
The one-loop bubble Feynman graph has two cubic vertices $V_{abc}$, two
internal propagator lines $P^{ab}=\delta^{ab}P(z_{1},z_{2})$, and an
external pair $(\cA^{c},\cA^{c'})$. Colour factor:
\[
f_{abc}f_{a'b'c'}\delta^{aa'}\delta^{bb'}\;=\;f_{abc}f^{ab}{}_{c'}\;=\;
\Cas\,\delta_{cc'}.
\]
Integrand: $\cA^{c}(z_{1})\cA^{c'}(z_{2})\cdot P(z_{1},z_{2})^{2}$
(wedge of two propagators). The $-1/(2\pi)^{3}$ normalisation comes from
the Bochner--Martinelli prefactor $2!/(2\pi i)^{3}$ and the bubble symmetry
factor $1/2$. Sign: the bubble contributes at $\hbar^{1}$ with a minus
sign from the BV Laplacian $\Delta_{\BV}$, which closes the bubble loop
by tracing over the two propagator endpoints.
\end{proof}

\begin{remark}[Why $\hdual$ and not $\Cas$?]\label{rem:hdual-vs-cas}
The relation $\Cas=2\hdual$ in the basis where the Killing form $\tr_{\mathrm{adj}}$
has normalisation $\tr(T^{a}T^{b})=\delta^{ab}$ is standard (Humphreys 1978
\S22.1, eq.~6). In 3d Chern--Simons the level shift $k\mapsto k+\hdual$
arises from a \emph{one-loop} computation with an explicit factor of $2$
absorbed into the definition of the level; in 6d hCS, the same Feynman-level
factor appears, but there is no discrete level parameter to shift
(companion wave12\_b4 \S4), and the coefficient $\Cas$ (or equivalently
$2\hdual$) appears as the coefficient of the anomaly class directly. The
$(2\pi)^{-3}$ is the dimensional factor from the $6$d volume element against
the cubic vertex.
\end{remark}

\begin{remark}[Alternative normalisations]\label{rem:alt-norms}
Some references absorb $2\pi i$-factors differently: the Costello 2013
``Supersymmetric and holomorphic field theories'' uses $(2\pi i)^{-n}$ for
$n$-complex-dimensional ambients; Costello--Gwilliam Vol.~II \S10.5 folds
the $(2\pi i)$-factors into the regularisation counterterms. The
cohomological content of the anomaly is regulator-independent; the explicit
coefficient in Proposition~\ref{prop:anomaly-poly} is fixed by the
Bochner--Martinelli prefactor.
\end{remark}

%% =====================================================================
\section{Cycle 4 --- $\chi_{\mathrm{top}}$ via Atiyah--Singer (K4, part 1)}
\label{sec:cycle4}
%% =====================================================================

\begin{attack}[Cycle 4 --- Why Euler characteristic, not signature or holomorphic Euler?]\label{atk:4}
The claim $\kanom\propto\chi_{\mathrm{top}}(X)$ uses the \emph{topological}
Euler number $\int_{X}c_{3}(TX)$, which is the alternating sum of Betti
numbers. Why not $\chi(\cO_{X})=\int_{X}\Td(TX)$ (Todd / holomorphic
Euler characteristic)? Why not $\sigma(X)$ (signature)? The essay does not
distinguish; we must.
\end{attack}

\begin{heal}[Cycle 4 --- $\chi_{\mathrm{top}}$ from Atiyah--Singer on $\End\gfrak$]\label{heal:4}
The one-loop anomaly lives in the index of the $\dbar$-Dirac operator
\emph{coupled to $\End\gfrak$}: the bubble closes the loop by tracing over
internal states which are propagator endpoints, each valued in
$\gfrak\otimes\gfrak\cong\End\gfrak$ (via the Killing form). The index
theorem gives
\begin{equation}\label{eq:AS-index}
\mathrm{ind}(\dbar_{\End\gfrak})\;=\;\int_{X}\ch(\End\gfrak)\cdot\Td(TX).
\end{equation}
For a trivial adjoint bundle on flat $X$ this would be
$\dim(\End\gfrak)\cdot\chi(\cO_{X})=\dim(\End\gfrak)$; but the anomaly
cocycle in \emph{local} BV cohomology detects the topological class,
which for a compact CY3 with $c_{1}(X)=0$ simplifies to
$\chi_{\mathrm{top}}(X)$.
\end{heal}

\begin{lemma}[Todd-vs-Euler on CY3]\label{lem:todd-vs-euler}
For a compact Calabi--Yau 3-fold $X$:
\[
\chi_{\mathrm{top}}(X)\;=\;\int_{X}c_{3}(TX),\qquad
\chi(\cO_{X})\;=\;\int_{X}\Td(TX)\;=\;\frac{1}{24}\int_{X}c_{1}c_{2}+\tfrac{1}{24}c_{3}
\;=\;0,
\]
the holomorphic Euler $\chi(\cO_{X})=0$ by Serre duality on CY3
($h^{0,0}=h^{0,3}=1$, $h^{0,1}=h^{0,2}=0$; total Hodge-supertrace
$1-0+0-1=0$). The topological Euler $\chi_{\mathrm{top}}=\int c_{3}$ is
non-zero in general ($\chi_{\mathrm{top}}(K3\times E)=0$ by K\"unneth,
$\chi_{\mathrm{top}}(\mathrm{quintic})=-200$, $\chi_{\mathrm{top}}(\PP^{2}_{\mathrm{local}})=3$
from McKay).
\end{lemma}

\begin{remark}[Two zeros for K3$\times$E, with different mechanisms]\label{rem:k3xE-zeros}
$\chi_{\mathrm{top}}(K3\times E)=\chi_{\mathrm{top}}(K3)\cdot\chi_{\mathrm{top}}(E)
=24\cdot0=0$, by the K\"unneth-multiplicative identity. $\chi(\cO_{K3\times E})
=\chi(\cO_{K3})\cdot\chi(\cO_{E})=2\cdot0=0$. Both zero, but via
\emph{different} K\"unneth factorisations: the topological Euler through
$\chi(E)=0$ (two tori eigenvalues cancel), the holomorphic Euler through
$\chi(\cO_{E})=0$ (Serre duality on an elliptic curve). In both cases,
the vanishing is a feature of $E$ as the elliptic fibre, not of the
Calabi--Yau structure of the threefold. CLAUDE.md cache entry: $\kcat(K3\times E)=0$
(total space, K\"unneth-multiplicative), NOT $\kcat(K3)=2$ (fibre).
\end{remark}

\begin{proposition}[Atiyah--Singer for the bubble on $X$]\label{prop:AS-for-bubble}
The one-loop anomaly integrand on a compact CY3 $X$ evaluates to
\[
\int_{X}\tr_{\gfrak\otimes\gfrak^{\ast}}\bigl[\ch(\End\gfrak)\wedge
\Td(TX)\bigr]_{(3,3)}\;=\;A(\gfrak)\cdot\chi_{\mathrm{top}}(X)
+(A(\gfrak)\text{-independent corrections}).
\]
For the trivial bundle (constant $\gfrak$-connection), $\ch(\End\gfrak)
=\dim(\End\gfrak)$ and $\Td(TX)|_{(3,3)}=c_{3}(TX)/24+c_{1}c_{2}/24$;
combined with the BV-local-cohomology projection onto the ghost-number-$+1$
class, only the $c_{3}$-component survives, yielding
$\chi_{\mathrm{top}}(X)=\int_{X}c_{3}(TX)$.
\end{proposition}

\begin{proof}
(Sketch.) The BV local-cohomology class is computed by the Gelfand--Fuks
calculation adapted to the holomorphic setting (Costello--Gwilliam 2017
Vol.~II Prop.~7.3.7); the surviving term is the Atiyah--Singer contribution
from the $\dbar$-Dirac on $\End\gfrak$. On a CY3, $c_{1}(TX)=0$, so
$\Td(TX)=1+c_{2}/12+c_{1}c_{2}/24+\ldots$ truncates at $c_{2}/12$ in
cohomological degree up to $(2,2)$; the $(3,3)$-component comes from
$c_{3}/24$. The Casimir trace $\Cas=2\hdual$ appears as the coefficient of
the fundamental-representation trace in the bubble, extracting $A(\gfrak)$.
The $c_{1}c_{2}$-term vanishes because $c_{1}=0$ on CY3. The net result
reduces to $A(\gfrak)\cdot\int_{X}c_{3}(TX)/\text{normalisation}$.
\end{proof}

%% =====================================================================
\section{Cycle 5 --- Multiplicative factorisation $A(\gfrak)\chi_{\mathrm{top}}(X)$ (K4, part 2)}
\label{sec:cycle5}
%% =====================================================================

\begin{attack}[Cycle 5 --- Why multiplicative factorisation?]\label{atk:5}
Proposition~\ref{prop:AS-for-bubble} gives a multiplicative decomposition
$\kanom(X,\gfrak)=A(\gfrak)\cdot\chi_{\mathrm{top}}(X)$. Why? What prevents
``$\gfrak$-dependent correction'' entering additively, or mixing with
non-topological data of $X$ (curvature tensor beyond characteristic classes)?
Is the factorisation structural or a feature of our one-loop computation
only?
\end{attack}

\begin{heal}[Cycle 5 --- Factorisation is Chern--Weil + index theorem]\label{heal:5}
The multiplicative factorisation is structural; it is the Chern--Weil
identification of the anomaly with a product of two independent
characteristic-class contributions.
\end{heal}

\begin{theorem}[Factorised one-loop anomaly]\label{thm:factorisation}
\ClaimStatusTheorem\ For any compact Calabi--Yau 3-fold $X$ and any
finite-dimensional Lie algebra $\gfrak$ with non-degenerate invariant
pairing, the one-loop BV anomaly of $\hCS(X,\gfrak)$ factorises as
\begin{equation}\label{eq:factorisation}
\kanom(X,\gfrak)\;=\;A(\gfrak)\cdot\chi_{\mathrm{top}}(X),
\end{equation}
with $A(\gfrak)=-\Cas/(2\pi)^{3}$ (Proposition~\ref{prop:anomaly-poly})
and $\chi_{\mathrm{top}}(X)=\int_{X}c_{3}(TX)$. On $\C^{3}$ (non-compact)
both factors vanish: $\chi_{\mathrm{top}}(\C^{3})=0$ (non-compact) matches
the Dolbeault dimension vanishing of wave12\_b4 Prop.~4.4.
\end{theorem}

\begin{proof}
The factorisation separates two independent pieces of data:
(i) the colour trace $\tr_{\End\gfrak}(\mathrm{bubble})=\Cas$, which
depends only on $\gfrak$ and is the coefficient of the Casimir in the
bubble colour factor; (ii) the spatial integral
$\int_{X}\mathrm{Integrand}(TX)$, which evaluates to the topological Euler
by Proposition~\ref{prop:AS-for-bubble}.

The separation is a Chern--Weil statement: the connection $\cA$ on the
trivial $\End\gfrak$-bundle over $X$ contributes a curvature $F_{\cA}$,
and the one-loop integral over connections involves
$\int_{X}\tr(F^{3})\wedge\mathrm{Ch.W.(TX)}$, but the cubic $\tr(F^{3})$
on a trivial bundle (flat $\cA$) reduces to a Casimir contribution times
the topological integral of the base. Since $c_{1}(TX)=0$, the Todd class
reduces to the curvature polynomial yielding $\int c_{3}$; since the
bubble closes with adjoint-representation trace, the Lie-theoretic
coefficient is $\Cas$. The two integrations are independent, forcing
multiplicativity.

The result is regulator-independent (Costello 2011 Thm.~8.6 ``independence
of gauge-fixing''; Costello--Gwilliam 2017 Vol.~II Thm.~5.4.1); the
heat-kernel regulator of Cycle~2 yields the explicit coefficient
$(-1/(2\pi)^{3})$ in $A(\gfrak)$.
\end{proof}

\begin{corollary}[Vanishing on $\C^{3}$]\label{cor:vanish-C3}
For $X=\C^{3}$, $\chi_{\mathrm{top}}(\C^{3})=0$ (contractible), so
$\kanom(\C^{3},\gfrak)=0$ for every $\gfrak$. This recovers
wave12\_b4 Theorem~4.5 from a dual perspective:
the Dolbeault-dimension-vanishing on $\C^{3}$ \emph{is} the
$\chi_{\mathrm{top}}=0$ statement on the non-compact ambient.
\end{corollary}

\begin{corollary}[Vanishing on $K3\times E$]\label{cor:vanish-k3xe}
For $X=K3\times E$, $\chi_{\mathrm{top}}(K3\times E)=24\cdot0=0$ by
K\"unneth, hence $\kanom(K3\times E,\gfrak)=0$. The 6d hCS on $K3\times E$
is \emph{one-loop anomaly-free for every $\gfrak$}. This is compatible
with but substantially stronger than the statement ``no BCOV anomaly on
$K3\times E$''.
\end{corollary}

\begin{corollary}[Non-vanishing on the quintic and on local $\PP^{2}$]\label{cor:nonvanish}
For $X=$ the quintic,
$\kanom(\mathrm{quintic},\gfrak)=-200\cdot A(\gfrak)$; for non-abelian
$\gfrak$ this is non-zero and obstructs quantisation of hCS on the quintic
without a non-trivial counterterm. For local $\PP^{2}$
($X=\mathrm{Tot}(K_{\PP^{2}})=\mathrm{Tot}(\cO_{\PP^{2}}(-3))$), the
topological Euler is $3$ (McKay), so $\kanom=3\cdot A(\gfrak)\neq0$ for
non-abelian. The CY-D stratification $\kch=3/2$ on local $\PP^{2}$
(CLAUDE.md essential constants) is the chiral-side avatar of this
one-loop BV statement.
\end{corollary}

\begin{remark}[Seven faces of $r_{\mathrm{CY}}$]\label{rem:seven-faces}
The factorisation $\kanom=A(\gfrak)\cdot\chi_{\mathrm{top}}(X)$ is the
one-loop BV face of $r_{\mathrm{CY}}$: the $\gfrak$-dependent factor $A(\gfrak)$
matches the ``chiral-algebra structure constant'' face; the
$\chi_{\mathrm{top}}$-dependent factor matches the ``CY3 topological'' face.
The multiplicative factorisation is the index-theorem shadow of a
categorical tensor product $(\text{Rep}\gfrak)\otimes(\mathrm{DCoh}(X))$ at
the level of anomaly classes.
\end{remark}

%% =====================================================================
\section{Comparison with companion notes}
\label{sec:comparison}
%% =====================================================================

\begin{itemize}
\item \texttt{wave12\_b4\_6d\_hCS\_BV\_BRST\_Feynman.tex} Proposition~4.4
proves the bubble on $\C^{3}$ vanishes by the Dolbeault-dimension
argument ($\Omega_{w}\wedge\Omega_{w}=0$, a property of non-compact
flat $\C^{3}$). The present note Theorem~\ref{thm:factorisation} establishes
this as the special case $\chi_{\mathrm{top}}(\C^{3})=0$ of the general
factorisation; the Dolbeault-dimension vanishing generalises to a
$\chi_{\mathrm{top}}$-dependent statement on compact $X$.

\item \texttt{wave12\_b3\_6d\_hCS\_E3\_gen\_rel.tex} Theorem~6.1 identifies
$\Obs^{q}_{(0)}\simeq U_{E_{3}}(\gfrak[2])$ at a point. The chiral
$E_{3}$-structure is independent of $\chi_{\mathrm{top}}$; the one-loop
anomaly of Theorem~\ref{thm:factorisation} is a \emph{global} obstruction
detecting $\chi_{\mathrm{top}}\cdot A(\gfrak)\neq0$ and is separate from
the local $E_{3}$-structure. Different $\gfrak$ on same $X$ give different
$E_{3}$-algebras at a point; the anomaly determines which of these lift
to a global factorisation algebra.

\item \texttt{wave12\_a13\_costello\_5\_routes\_6d\_hCS.tex} Attack~2
identifies the Costello--Li perturbative closure. Theorem~\ref{thm:factorisation}
is the compact-ambient refinement: on $\C^{3}$ every $\gfrak$ closes
perturbatively; on compact CY3 with $\chi_{\mathrm{top}}\neq0$, non-abelian
$\gfrak$ requires a counterterm (or fails to close).

\item \texttt{wave14\_i1\_hCS\_nonabelian\_anomaly.tex} \S2 proves the
Costello obstruction cocycle on $\C^{3}$ vanishes. The present note extends
the statement to compact $X$: on compact $X$ the cocycle is
$A(\gfrak)\cdot\chi_{\mathrm{top}}(X)$, non-zero for non-abelian $\gfrak$
and $\chi_{\mathrm{top}}\neq0$.

\item \texttt{wave13\_b3\_6d\_hCS\_deep.tex} and
\texttt{platonic\_synthesis\_waves\_11\_through\_16.tex} are synthesis-level;
the present note is a first-principles upstream statement.
\end{itemize}

%% =====================================================================
\section{Cache entries}
\label{sec:cache}
%% =====================================================================

Five AP-CY candidates for \texttt{antipatterns\_catalogue.md}. Each follows
the format (confusion, correction, primary source).

\begin{enumerate}
\item \textbf{Bochner--Martinelli as class representative.}
Confusion: ``BM is the unique propagator on $\C^{3}$.'' Correction: BM is
the class representative in the affine gauge-fixing space, tied to the flat
Hermitian metric; $Q$-exact deformations give homotopic propagators with
same cohomological content. Primary: Costello 2011 Thm.~A.7.1;
Proposition~\ref{prop:BM-derivation} above.

\item \textbf{Heat-kernel regulator degenerates to BM.} Confusion:
``$P_{\varepsilon,L}$ and $P_{\mathrm{BM}}$ are different propagators.''
Correction: $P_{\varepsilon,L}\to P_{\mathrm{BM}}$ as $\varepsilon\to0$,
$L\to\infty$ off the diagonal; same cohomology class throughout the
renormalisation flow. Primary: Proposition~\ref{prop:propagator-identity}
above; Costello--Gwilliam 2017 Vol.~II \S8.

\item \textbf{$A(\gfrak)$ is Casimir, not dual Coxeter (separately).}
Confusion: ``$A(\gfrak)\propto\hdual$ analogous to 3d CS.'' Correction:
$A(\gfrak)=-\Cas/(2\pi)^{3}$ with $\Cas=2\hdual$ in normalised basis; the
relation $\Cas=2\hdual$ is automatic, the overall sign and $(2\pi)^{-3}$
come from Bochner--Martinelli + bubble symmetry factor. Primary:
Proposition~\ref{prop:anomaly-poly} above.

\item \textbf{$\chi_{\mathrm{top}}$ vs $\chi(\cO_{X})$.} Confusion:
``The anomaly depends on the holomorphic Euler $\chi(\cO_{X})$.'' Correction:
It depends on the topological Euler $\chi_{\mathrm{top}}(X)=\int_{X}c_{3}$;
$\chi(\cO_{X})=0$ on every CY3 by Serre duality, hence cannot carry the
anomaly content. Primary: Lemma~\ref{lem:todd-vs-euler} above.

\item \textbf{Multiplicative factorisation is Chern--Weil.} Confusion:
``$\gfrak$-dependent and $X$-dependent contributions could mix.'' Correction:
Chern--Weil on the trivial $\End\gfrak$-bundle separates the colour trace
from the base-space integral; the index-theorem calculation forces
multiplicativity. Primary: Theorem~\ref{thm:factorisation} above;
Atiyah--Singer 1968 III.
\end{enumerate}

%% =====================================================================
\section*{Epistemic scope}
%% =====================================================================

\ClaimStatusTheorem\
Propositions~\ref{prop:BM-derivation}, \ref{prop:propagator-identity},
\ref{prop:anomaly-poly}, \ref{prop:AS-for-bubble},
Theorem~\ref{thm:factorisation}, Corollaries~\ref{cor:equiv-class},
\ref{cor:regulator-compat}, \ref{cor:vanish-C3},
\ref{cor:vanish-k3xe} follow from standard Bochner--Martinelli theory,
Costello--Gwilliam renormalisation, and Atiyah--Singer on the
$\dbar$-Dirac operator coupled to $\End\gfrak$. The first-principles
assembly into the factorised form $A(\gfrak)\chi_{\mathrm{top}}(X)$ is the
contribution of this note.

\ClaimStatusHeuristic\
The informal ``dimensional reduction'' 3d$\to$6d analogy suggesting
$A(\gfrak)\propto\hdual$ without the $2$-factor or $(2\pi)^{-3}$ is a
conventional shorthand; the first-principles computation
(Proposition~\ref{prop:anomaly-poly}) fixes the correct
normalisation as $-\Cas/(2\pi)^{3}$.

\ClaimStatusConjectured\
The identification of the index-theorem shadow with the categorical
tensor product $(\text{Rep}\gfrak)\otimes(\mathrm{DCoh}(X))$ at the level
of anomaly classes (Remark~\ref{rem:seven-faces}) is a programmatic
statement about the seven faces of $r_{\mathrm{CY}}$, not a theorem; the
rigorous two-lane statement (chain-level index identity, $(\infty,1)$-categorical
tensor-product statement) is work-in-progress in the companion
\texttt{chapters/theory/cy\_to\_chiral.tex}.

\ClaimStatusDefinition\
BV data, propagator family, heat-kernel regulator, anomaly polynomial
as defined above are standard (Costello 2011;
Costello--Gwilliam 2017 Vol.~II; Costello--Li 2015; Costello 2013
``Supersymmetric and holomorphic field theories'').

\end{document}
